\section*{Discussion}

On average, a cortical region's response to local optogenetic inhibition is approximately linear and time-invariant; single trials depart from it together with the state-dependent ongoing activity they ride on. A controller that treated the average as the plant and single-trial departures as a disturbance held a somatosensory site nearer its target than fixed stimulation in all 15 sessions, and partly rejected the component of its activity predicted from the opposite hemisphere. What feedback left was informative about state: error that accompanied movement was attenuated, whereas slow, large fluctuations were reduced but not held. Holding activity at a chosen value tightens causal tests of mesoscale dynamics and turns state-dependent variability from a confound averaged away into a disturbance partly removed on each trial.

Linearity was established for the average, not for single trials. Suppression grew monotonically with laser power and a low-order linear model fitted the trial average (Figure~\ref{fig:figure2}B--D) but transferred inconsistently between sessions (Figure~\ref{fig:figure2}E). Linear models likewise describe resting human brain activity at least as well as nonlinear ones, which their authors attribute to averaging, noise and limited data~\citep{Nozari2024}; trial averaging is one more such step. We treat the site-to-readout system as nonlinear, non-stationary and driven by unobserved inputs, and design on its trial average, because a low-order linear model makes stability, gain limits and bandwidth computable~\citep{astrom2008feedback}. Monotonicity gives each feasible reference one steady-state light drive, which is what integral action needs. Low-gain integral control of an exponentially stable linear plant reaches it through a saturating actuator and input delay~\citep{logemann2000absolute}, although no single result covers our sampled proportional--integral loop (Remark~\ref{rem:lowgain}). Session-to-session variability therefore limits the gains, not the existence of a stabilizing controller of fixed sign~\citep{morari1985robust}. That sign was measured over our powers, not assumed (Figure~\ref{fig:figure2}B; Figure~\ref{fig:figure3}B); paradoxical responses of inhibition-stabilized cortex concern interneuron rates and were absent on average with viral expression in parvalbumin interneurons~\citep{sanzeni2020inhibition}. State dependence is present in the unstimulated cortex and appears in stimulus responses: single-trial responses and stimulus-free activity co-varied with movement and absolute 2--4\,Hz power, so their ratio barely changed (Figure~\ref{fig:figure2}G; Figure~\ref{fig:state_quartile}). This is consistent with responses riding on preceding ongoing activity, as in anesthetized cat visual cortex~\citep{Arieli1996} and awake mouse dorsal cortex~\citep{Shimaoka2019}. A change in the stimulus effect itself is not excluded: responses to fixed optogenetic drive of excitatory neurons in mouse barrel cortex vary with prestimulus membrane potential~\citep{Mateo2011}, and variability in macaque visual cortex is partly multiplicative~\citep{Goris2014}. State that adds an offset is a disturbance that feedback subtracts; state that changes gain is a model error that feedback must tolerate. Input--output models fitted per state quartile would tell them apart.

Feedback partly cancels the activity that the opposite hemisphere predicts at the controlled site. Relative to this predicted disturbance, feedback lowered residual energy in all eleven sessions, mostly by removing the standing offset left by the open loop's feedforward calibration; fluctuations fell less (Figure~\ref{fig:figure4}G; Table~\ref{tab:stat_fig4g}). An integrator drives a constant offset toward zero within its clamp~\citep{francis1976internal}, whereas the loop attenuates other disturbances only where its gain is large~\citep{astrom2008feedback}, so the smaller fall in fluctuation better measures cancelled shared drive. Interhemispheric restoration of silenced activity in mouse anterolateral motor cortex~\citep{li2016robust} suggests that one side opposes an open-loop perturbation of the other. The contralateral split is a linear projection, not a causal decomposition: in awake mice, optogenetically evoked interhemispheric connectivity, read out hemodynamically, is weaker than resting-state correlations where those are strong~\citep{bauer2018effective}, so the shared component may reflect common subcortical, neuromodulatory or movement-related drive. Contralateral silencing would distinguish these.

Feedback separates ongoing cortical variability into components it can and cannot reach. Error that accompanied movement was attenuated: the closed-to-open-loop error ratio fell 9\% per SD of motion, because open-loop error rose (predictability) while closed-loop error stayed flat (regularizability; Figure~\ref{fig:figure4}D). Impulse responses, by contrast, scattered least at high motion (Figure~\ref{fig:state_quartile}A), perhaps because slow fluctuations shrink in desynchronized states~\citep{McGinley2015waking}; the impulse sessions and state window differ from the controller's. Slow, large-amplitude 1--4\,Hz fluctuations were reduced but not held: closed-loop error rose with this power as much as open-loop error, so the error ratio, a broadband analogue of the sensitivity function~\citep{astrom2008feedback}, grew 8\% per SD, whereas a large initial deviation lost the same fraction as a small one (Table~\ref{tab:fig4_state_stats}). Loop timing is one limit. An effective delay of 37--61\,ms (loop delay at most 47\,ms and about half that on average, Figure~\ref{fig:figure1}F, plus the 14\,ms half-sample hold) caps, by the crossover criterion $\omega_c\tau \lesssim 1$\,rad, any controller's bandwidth at about 2.6--4\,Hz~\citep{skogestad2005multivariable}; cortical lag (about 60 degrees at 1\,Hz; Figure~\ref{fig:figure5}H) probably lowers the unmeasured crossover further. Alternatively, large slow excursions may saturate the one-sided command, the band-power window spans stimulation, which the loop shapes, and motion and slow power are anticorrelated~\citep{McGinley2015waking}. The one-sided actuator also limits reachability: the two sessions excluded from the disturbance analysis for a standing offset stayed above the reference (slow integration or the integral clamp could also explain this), and in the sinusoidal replay about half the squared error the model-predictive controller left fell in rising half-cycles, when the laser was mostly off. These slopes are associations; a closed- to open-loop error spectrum and time at the ceiling would separate delay from saturation.

Drive from the rest of cortex is observable but, at short leads, largely not forecastable. The contralateral hemisphere predicted much of the site's spontaneous activity (Figure~\ref{fig:figure4}F), yet on one session's held-out open-loop trials, 100 cortical sites added little to a 200\,ms forecast of the site from its own past and the laser (Figure~\ref{fig:mpc_direct}H; Figure~\ref{fig:supp_varx}E), though a linear model and unrecorded subcortical signals may understate what is forecastable. Latent-dynamics models infer per-trial inputs from the whole trial~\citep{pandarinath2018lfads}, so their fit over-approximates what is knowable in real time. The best of a dozen univariate, stimulus-blind forecasters of spontaneous mouse wide-field activity stays informative for about 1.5\,s~\citep{lu2025benchmarking}, and the authors of a larval zebrafish whole-brain benchmark expect an irreducible error floor~\citep{lueckmann2025zapbench}; neither models an actuator input. Preview, which predicts only the reference, cancelled the cortical phase lag that feedback left (Figure~\ref{fig:figure5}H). In a replay of one session, a model-predictive controller, which anticipates the predictable disturbance and treats the one-sided laser as a constraint~\citep{mayne2000constrained}, lowered the median per-trial RMSE that feedback left over 1--3\,s by about 40\%, and by about 80\% given the future residuals (Figure~\ref{fig:mpc_direct}C--F). What it left was mostly activity its model could not predict. Model-based optogenetic feedback in awake mouse thalamus instead adapted its estimator online to absorb model mismatch~\citep{bolus2021state}. The replayed cortex shares the controller's model class and responds to the laser identically in every state, so these values bound rather than measure a live controller.

The loop fed back on uncorrected fluorescence, because hemodynamic correction cannot run within one control interval, so closed-loop measures describe the regulated signal rather than purely neural activity~\citep{valley2020separation}. Inhibition spreads toward the adjacent readout (Figure~\ref{fig:spatial_spread}A,B;~\citealp{li2019spatiotemporal}), so the loop holds a readout kernel, not a whole area, and partly regulates that spread. The sample is small: four impulse sessions from three mice, separate from the controller sessions; 15 controller sessions from four mice, eleven of them in the disturbance analysis; and three sinusoidal sessions from one mouse. All inhibited one somatosensory site in awake, head-fixed mice, leaving excitation, other areas and task performance untested. Interleaving open- and closed-loop trials guards against drift in brain state between conditions but cannot remove it. Feedback gains were tuned empirically and the low-gain guarantee assumes a fixed plant, whereas brain state changes over seconds, so stability was shown empirically: no session showed sustained oscillation. There was no sham, opsin-free light-only, random-timing or proportional-only condition, so opsin-independent light effects would enter the identified plant, and because proportional action also shrinks an offset, the role of integral action in removing it is inferred, not isolated.

With a site held at a chosen level, its influence on dorsal cortex, striatum~\citep{Peters2021} or behavior~\citep{ZatkaHaas2021} can be studied against a set value. As with the voltage clamp and its network analogue, the optoclamp~\citep{newman2015optogenetic}, the holding command should track the drive the loop opposes, a prediction we did not test. A live predictive controller should keep integral action~\citep{pannocchia2003disturbance}, and lower latency and an excitatory actuator would widen what feedback can reach.

Cortical activity that differs on every trial is regular enough on average for feedback to steer it toward a target, and what feedback cannot remove maps which components of ongoing variability lie within its reach.
